\chapter{Stochastik}

\section{Bernoulli-Experiment}

\startbuffer[fig:stochastik:bernoulli-kette:baum]
\switchtobodyfont[8pt]
\startMPpage

u := 0.6cm;

color lightgreen, lightred;
lightgreen := (1, 0.8, 0.8);
lightred := (0.8, 1, 1);

% Define colors, if desired
pickup pencircle scaled 1pt;

vardef draw_node(expr x, y, myLable, myColor) =
	draw fullcircle scaled (0.6 * u) shifted (x * u, y * u);
	fill fullcircle scaled (0.6 * u) shifted (x * u, y * u) withcolor myColor;
	label(myLable, (x * u,y * u));
enddef;

% Function to draw the tree recursively
% vardef draw_tree(expr x, y, depth, text mylable) =
vardef draw_tree(expr x, y, depth, layer, myHeight, depthLength, myLable, myColor, ergebnisString) =
	save left_endpoint, right_endpoint;
  % Base case: if depth is 0, draw a leaf node (circle)
  if depth = 0:
	draw_node(x, y, myLable, myColor);
	label(ergebnisString, ((x + 1) * u, y * u));
  else:

	numeric myOffset;
	myOffset = ((myHeight / (2**layer)));

    pair left_endpoint, right_endpoint;
    left_endpoint = (x + depthLength, y - (myOffset / 2));
    right_endpoint = (x + depthLength, y + (myOffset / 2));
    
	draw (x * u, y * u) -- right_endpoint * u;
	label.top("p", .5[(x * u, y * u), right_endpoint * u]);

	draw (x * u, y * u) -- left_endpoint * u;
	label.bot("1-p", .5[(x * u, y * u), left_endpoint * u]);

	draw_node(x, y, myLable, myColor);

	draw_tree(xpart right_endpoint, ypart right_endpoint, depth - 1, layer + 1, myHeight, depthLength, "1", lightgreen, ergebnisString & "1");
	draw_tree(xpart left_endpoint, ypart left_endpoint, depth - 1, layer + 1, myHeight, depthLength, "0", lightred, ergebnisString & "0");
  fi
enddef;

% draw_tree(0, 0, 2, "");
% x, y, depth, layer, myHeight, depthLength, myLable, myColor
draw_tree(0, 0, 3, 1, 8, 2, "", white, "");

label("1. Wurf",	((0 * 2 + 0.75) * u, 5 * u));
label("2. Wurf",	((1 * 2 +	0.75) * u, 5 * u));
label("3. Wurf",	((2 * 2 +	0.75) * u, 5 * u));
label("Ergebnis",	((3 * 2	+ 0.75) * u, 5 * u));

\stopMPpage
\stopbuffer

\margintext{
\placefigure[here][fig:stochastik:bernoulli-kette:baum]{Bernoulli-Experiment mit drei Wiederholungen}{
\typesetbuffer[fig:stochastik:bernoulli-kette:baum]
}}

Bei einem Bernoulli-Experiment gibt es nur zwei Ereignisse\sidenote{Diese
müssen nicht gleich wahrscheinlich sein.}. Zum Beispiel beim würfeln, ob ein
$6$ gewürfelt wird, oder nicht. Das Ereignis $6$ gewürfelt wird in einem
Baumdiagram mit $1$ dargestellt, das Gegenereignis mit $0$.

Das Baumdiagram in der \in{Abbildung}[fig:stochastik:bernoulli-kette:baum]
stellt den Würfelversuch mit drei Wiederholungen dar. Das Ergebnis "111" wird
als "in allen drei Würfen eine $6$" gelesen. 

Die Wahrscheinlichkleit eine $6$ zu würfeln liegt bei:

\startformula

p = \frac{1}{6}

\stopformula

Dadruch ergibt sich auch die Wahrscheinlichkeit für das Gegenereignis:

\startformula

1 - p = \frac{5}{6}

\stopformula


\startdefinition[title={Bernoulli-Experiment}]
Ein {\bf Zufallsexperiment} heißt {\bf Bernoulli-Experiment}, wenn es genau zwei
Ergebnisse hat. Ein {\bf Bernoulli-Kette} besteht aus mehreren voneinander
unabhängigen Durchführungen eines Bernoulli-Experiment. Die Anzahl der Durchführungen nennt man {\bf Länge} $n$ der Bernoulli-Kette, die {\bf Treffwahrscheinlichkeit} wird mit $p$ angegeben.
\stopdefinition

\margintext{Ab Seite: 273}
\section{Binominalkoeffizient}


\stoptabulate
